\documentclass[11pt,a4paper]{article}
\usepackage[margin=1in]{geometry}
\usepackage{times}
\usepackage{latexsym}
\usepackage{booktabs}
\usepackage{amsmath}
\usepackage{amssymb}
\usepackage{graphicx}
\usepackage{multirow}
\usepackage{xcolor}
\usepackage{subcaption}
\usepackage[T1]{fontenc}
\usepackage[utf8]{inputenc}
\usepackage{hyperref}
\usepackage{natbib}

\title{Structural Features Dominate Zero-Shot Cross-Agent Stagnation Detection: \\A Gold-Validated Empirical Study}

\author{Anonymous}

\begin{document}
\maketitle

\begin{abstract}
LLM-based coding agents frequently enter \emph{progress stagnation}---executing seemingly reasonable but collectively unproductive action sequences.
We present a systematic empirical study of stagnation detection through weak supervision, comparing structural features, neural probes, and LLM-as-judge across two agents (SWE-Agent, OpenHands) totaling 194k trajectory windows.
Our central finding, validated on 650 human-annotated windows ($\kappa$=0.82), is that 30 lightweight structural features achieve 0.948 gold-validated macro-F1, outperforming a GRU probe by 8.7pp and an LLM judge by 25.9pp in zero-shot cross-agent transfer---with zero target-domain data and sub-5ms inference.
Even after removing all 14 features overlapping with labeling function logic, the remaining 16 non-overlap features achieve 0.890 gold F1, confirming that structural detection captures genuine stagnation signals beyond circular label artifacts.
Cross-agent transfer proves robust to instance isolation ($\Delta$=$-$0.62pp, $p$=0.43), ruling out data leakage.
We further show that few-shot prompting improves LLM judges only marginally (0.689$\to$0.723), leaving a 22.5pp gap to structural methods---the efficiency deficit is fundamental, not prompt-engineering addressable.
Our results establish structural features as the default choice for zero-data agent monitoring deployment and provide gold-validated practical guidance for the agent monitoring community.
\end{abstract}

%------------------------------------------------------------------
\section{Introduction}
\label{sec:intro}
%------------------------------------------------------------------

LLM-based coding agents such as SWE-Agent \cite{yang2024swe} and OpenHands \cite{wang2024openhands} can autonomously navigate codebases, write patches, and run tests.
Yet a persistent failure mode undermines their reliability: \emph{progress stagnation}, where an agent executes a sequence of individually plausible actions that collectively fail to advance toward the goal \cite{jimenez2024swebench}.
Unlike simple repetition loops, stagnation can manifest as semantically diverse but ineffective exploration---editing different files, running varied commands, yet making no meaningful progress.
Recent work on SlopCodeBench \cite{slopcodebench2026} documents that structural erosion occurs in 77\% of long-horizon coding trajectories, with verbosity increasing in 75.5\%, confirming that stagnation is pervasive rather than anecdotal.

Detecting stagnation is critical for practical deployment: timely detection enables intervention strategies (e.g., backtracking, re-planning, or human escalation) that can recover otherwise wasted computation.
Existing approaches span cycle detectors \cite{zhuge2024cdcs,specra2025,trajad2025} that miss semantically diverse stagnation, step-level process reward models \cite{xi2026agentprm,prm2025framework} that require outcome labels and RL pipelines, and self-reflection methods \cite{shinn2023reflexion,yao2023tree} that rely on the stuck model's own judgment.

We frame stagnation detection as a \emph{window-level binary classification} task: given a sliding window of $k$ consecutive agent steps, predict whether the window exhibits progress or stagnation.
Labels are generated via a weak supervision pipeline with three labeling functions (LFs)---command repetition, file-diff frequency, and outcome backpropagation---aggregated by majority vote.
We systematically compare three detection paradigms: (1)~\textbf{structural features}: 30 hand-crafted behavioral statistics (command repetition ratios, output novelty, action entropy, etc.); (2)~a \textbf{GRU probe} over frozen LLM embeddings (Qwen3-Embedding-0.6B); and (3)~an \textbf{LLM-as-judge} baseline (Claude Sonnet 4.6).

Our study yields three findings:

\begin{enumerate}
\item \textbf{Structural features dominate zero-shot transfer.}
On 650 gold-annotated windows, structural features (Gradient Boosting) achieve 0.948 macro-F1, outperforming the GRU probe (0.861, $-$8.7pp) and LLM-as-judge (0.689 zero-shot, 0.723 few-shot).
Cross-agent transfer from SWE-Agent to OpenHands reaches 0.967 on LF labels and remains robust under instance isolation ($\Delta$=$-$0.62pp).

\item \textbf{The LLM efficiency gap is fundamental.}
Even with 5-shot exemplars and chain-of-thought prompting, the LLM judge improves only 3.4pp over zero-shot, leaving a 22.5pp gap to structural features---at 640$\times$ higher latency and \$18.92 per 650 windows versus \$0.

\item \textbf{Structural signal extends beyond labeling function logic.}
An overlap ablation removing all 14 LF-correlated features retains 0.890 gold F1 with the remaining 16 features---still above the probe (0.861) and far above the LLM judge---confirming that structural detection captures dimensions of stagnation that labeling functions do not encode.
\end{enumerate}

%------------------------------------------------------------------
\section{Related Work}
\label{sec:related}
%------------------------------------------------------------------

\paragraph{Agent Failure Detection.}
Cycle detectors identify repetitive patterns \cite{zhuge2024cdcs,specra2025,trajad2025} but miss semantically diverse stagnation.
Process reward models \cite{xi2026agentprm,prm2025framework,toolprmbench2026,trajreward2026} score individual steps using outcome-derived training signals, requiring labeled trajectories and RL infrastructure.
Self-reflection approaches \cite{shinn2023reflexion,yao2023tree} and value-function methods \cite{le2022coderl} rely on the agent's own assessment or reward shaping.
Our approach differs: window-level detection via weak supervision requires no outcome labels, no RL pipeline, and operates as an external monitor rather than an agent component.

\paragraph{Agent Monitoring.}
Recent work establishes the need for sequence-level monitoring beyond single-step analysis \cite{monitorharm2025}.
Weak-to-strong monitoring \cite{weakstrong2025} studies safety-focused oversight of capable agents.
SlopCodeBench \cite{slopcodebench2026} benchmarks agent degradation in long-horizon coding tasks, finding structural erosion in 77\% of trajectories.
Our work provides the first systematic comparison of detection paradigms (structural, neural, LLM) for progress monitoring.

\paragraph{Weak Supervision.}
Data programming \cite{ratner2016data,ratner2020snorkel} aggregates noisy labeling functions into probabilistic labels.
While typically applied to i.i.d.\ data, agent trajectories exhibit sequential dependencies that violate this assumption.
We document these temporal dynamics and test whether temporal label models can exploit them (finding they cannot, in our regime).

%------------------------------------------------------------------
\section{Method}
\label{sec:method}
%------------------------------------------------------------------

\subsection{Task Formulation}

Given an agent trajectory $\tau = (a_1, o_1, a_2, o_2, \ldots, a_T, o_T)$ of actions $a_t$ and observations $o_t$, we extract sliding windows of $k=5$ consecutive steps.
Each window $w_i = (a_i, o_i, \ldots, a_{i+k-1}, o_{i+k-1})$ is labeled as \textsc{stagnation} or \textsc{progress}.

\subsection{Weak Supervision Pipeline}

Three labeling functions vote on each window:

\begin{itemize}
\item \textbf{Command Repetition (cmd\_rep)}: Flags windows with high command repetition ratio. Operates as a high-precision, low-recall detector (precision 97.9\%, FP rate 0.15\%).
\item \textbf{File Diff (file\_diff)}: Flags windows where file modifications are absent or minimal, indicating the agent is not producing tangible code changes.
\item \textbf{Outcome Backpropagation (OB)}: Propagates terminal failure signals backward through the trajectory as distant supervision.
\end{itemize}

Labels are aggregated by majority vote (MV) with a confidence threshold: windows receiving $\geq$2 concordant votes form the training set.
On SWE-Agent, the overall stagnation prevalence is 28.5\%.

\subsection{Detection Methods}

\paragraph{Structural Features.}
We extract 30 features per window spanning five categories: command patterns (repetition ratio, unique ratio, bigram repeats, entropy), file change statistics (diff ratio, lines changed, unique files), output characteristics (length mean/std, novelty score), temporal signals (step duration, token count trend), and behavioral diversity (action type entropy, tool switch rate, directory changes, edit revert ratio).
Features are fed to Random Forest (RF) and Gradient Boosting (GB) classifiers.

\paragraph{GRU Probe.}
Each step is embedded using Qwen3-Embedding-0.6B (1024-dim).
A bidirectional GRU with attention pooling maps the window's embedding sequence to a binary prediction.
The probe is encoder-agnostic: BGE, E5, and GTE embeddings all achieve $>$0.90 in-domain F1.

\paragraph{LLM-as-Judge.}
Claude Sonnet 4.6 classifies each window via zero-shot or 5-shot + chain-of-thought prompting, given the raw action--observation text.

\subsection{Gold Validation Set}

To address circular evaluation concerns (structural features can reverse-engineer LF labels at 0.995 F1), we construct a gold set of 650 windows annotated by 3 human annotators using majority vote.
Inter-annotator agreement is substantial: Fleiss' $\kappa$=0.82, mean pairwise Cohen's $\kappa$=0.83, with 92.0\% agreement rate.
Gold stagnation prevalence is 29.1\%, closely matching the LF-derived distribution.

%------------------------------------------------------------------
\section{Experiments}
\label{sec:experiments}
%------------------------------------------------------------------

\subsection{Gold-Validated Method Comparison}
\label{sec:gold}

Table~\ref{tab:gold} presents all methods evaluated on the 650 gold-annotated windows.
Structural GB achieves 0.948 gold macro-F1 (95\% CI: [0.930, 0.965]), confirming that its near-perfect LF-eval performance (0.995) is not merely circular: the 4.7pp drop reflects the removal of circular LF-reconstruction boost, while the remaining 0.948 represents genuine stagnation detection validated against human judgment.

The GRU probe achieves 0.861 gold F1 (CI: [0.833, 0.887]), 8.7pp below structural features.
The LLM judge scores 0.689 zero-shot, 25.9pp below structural.

\begin{table}[t]
\centering
\small
\begin{tabular}{lccccc}
\toprule
\textbf{Method} & \textbf{Gold F1} & \textbf{F1-stag} & \textbf{F1-prog} & \textbf{LF F1} & \textbf{$\Delta$} \\
\midrule
Structural GB   & \textbf{0.948} & 0.925 & 0.972 & 0.995 & $-$4.7 \\
Structural RF   & 0.944 & 0.919 & 0.969 & 0.995 & $-$5.1 \\
GRU Probe       & 0.861 & 0.803 & 0.919 & 0.906 & $-$4.4 \\
LLM zero-shot   & 0.689 & 0.620 & 0.758 & 0.663 & +2.7 \\
TWV Probe       & 0.855 & 0.794 & 0.916 & 0.900 & $-$4.5 \\
\bottomrule
\end{tabular}
\caption{Gold-validated method comparison ($N$=650). $\Delta$ = Gold$-$LF F1 (pp). Structural GB dominates with 0.948 gold macro-F1, an 8.7pp advantage over the probe and 25.9pp over the LLM judge.}
\label{tab:gold}
\end{table}

On the 52 disagreement windows where gold majority vote diverges from the LF pipeline, structural features achieve 56\% accuracy versus 38\% for the probe, suggesting structural features capture genuine stagnation patterns even on boundary cases where human annotators themselves showed uncertainty.

\begin{table}[t]
\centering
\small
\begin{tabular}{lc}
\toprule
\textbf{Metric} & \textbf{Value} \\
\midrule
Fleiss' $\kappa$ (3 annotators)    & 0.82 \\
Mean pairwise Cohen's $\kappa$     & 0.83 \\
Gold-MV agreement rate             & 92.0\% \\
Disagreement windows               & 52/650 (8.0\%) \\
Gold stagnation prevalence         & 29.1\% \\
\bottomrule
\end{tabular}
\caption{Inter-annotator agreement statistics for the gold validation set. Fleiss' $\kappa$=0.82 indicates substantial agreement, confirming stagnation is a reliably annotatable concept at the window level.}
\label{tab:iaa}
\end{table}

\subsection{LLM-as-Judge: Zero-Shot vs.\ Few-Shot}
\label{sec:llm}

A natural reviewer concern is whether the LLM baseline is unfairly weak.
We evaluate a strengthened configuration: 5-shot exemplars with chain-of-thought prompting (Table~\ref{tab:llm}).
Few-shot + CoT improves macro-F1 from 0.689 to 0.723 (+3.4pp), but leaves a 22.5pp gap to structural features.
The improvement is modest because stagnation detection requires tracking goal-directed coherence across a window---a capability that benefits from structured feature extraction rather than longer reasoning chains.

\begin{table}[t]
\centering
\small
\begin{tabular}{lcccc}
\toprule
\textbf{Config} & \textbf{Gold F1} & \textbf{F1-stag} & \textbf{Cost} & \textbf{Latency} \\
\midrule
Zero-shot       & 0.689 & 0.620 & \$11.66$^\dagger$ & $\sim$2.1s \\
5-shot + CoT    & 0.723 & 0.671 & \$18.92 & $\sim$3.2s \\
\midrule
Structural GB   & \textbf{0.948} & \textbf{0.925} & \$0 & $<$5ms \\
GRU Probe       & 0.861 & 0.803 & \$0 & $\sim$12ms \\
\bottomrule
\end{tabular}
\caption{LLM-as-judge variants on gold set ($N$=650). Even optimized prompting leaves a 22.5pp gap to structural features at 640$\times$ higher latency and non-trivial API cost. $^\dagger$Extrapolated.}
\label{tab:llm}
\end{table}

Error analysis reveals two dominant failure modes of the LLM judge: false negatives (60\% of errors) where diverse-action stagnation appears productive to the LLM in a single reading pass, and false positives (40\%) where repetitive but effective debugging loops are misjudged as stagnation (Table~\ref{tab:llm-errors}).

\begin{table}[t]
\centering
\small
\begin{tabular}{lccp{4.2cm}}
\toprule
\textbf{Error} & \textbf{Count} & \textbf{\%} & \textbf{Pattern} \\
\midrule
FP & 72  & 40\% & Repetitive debugging loops that are actually effective \\
FN & 108 & 60\% & Diverse-action stagnation looking productive \\
\bottomrule
\end{tabular}
\caption{Error analysis of LLM few-shot judge on gold set. The dominant failure (60\%) is missing diverse stagnation---precisely the gap structural features fill.}
\label{tab:llm-errors}
\end{table}

\subsection{Overlap Ablation: Beyond Circular Evaluation}
\label{sec:ablation}

Structural features could achieve high F1 simply by reverse-engineering labeling function logic.
To quantify this concern, we partition the 30 features into 14 that directly overlap with LF heuristics (e.g., \texttt{cmd\_rep\_ratio}, \texttt{file\_diff\_count}) and 16 non-overlap features (e.g., \texttt{output\_novelty\_score}, \texttt{action\_type\_entropy}, \texttt{edit\_revert\_ratio}).

Table~\ref{tab:ablation} shows that non-overlap features alone achieve 0.890 gold F1---a 5.8pp drop from the full set but still above the GRU probe (0.861, +2.9pp) and far above the LLM judge (0.723, +16.7pp).
The top contributing non-overlap features are output novelty (importance 0.183), action type entropy (0.156), and edit revert ratio (0.134)---behavioral signatures that labeling functions do not encode.

\begin{table}[t]
\centering
\small
\begin{tabular}{lccc}
\toprule
\textbf{Feature Set} & \textbf{$N$} & \textbf{Gold F1} & \textbf{$\Delta$ vs Full} \\
\midrule
Full structural       & 30     & \textbf{0.948} & --- \\
Non-overlap only      & 16     & 0.890 & $-$5.8pp \\
LF-overlap only       & 14     & 0.929 & $-$2.0pp \\
GRU Probe (ref)       & 1024-d & 0.861 & $-$8.7pp \\
\bottomrule
\end{tabular}
\caption{Overlap ablation on gold set. Removing all 14 LF-overlap features retains 0.890 gold F1---above the probe---confirming structural detection captures signal dimensions beyond LF logic.}
\label{tab:ablation}
\end{table}

\begin{table}[t]
\centering
\small
\begin{tabular}{lcp{4.5cm}}
\toprule
\textbf{Feature} & \textbf{Imp.} & \textbf{Interpretation} \\
\midrule
output\_novelty       & 0.183 & Low novelty $\to$ stagnation \\
action\_type\_entropy & 0.156 & Low diversity $\to$ stagnation \\
edit\_revert\_ratio   & 0.134 & High revert rate $\to$ ineffective edits \\
token\_count\_trend   & 0.098 & Declining tokens $\to$ disengagement \\
tool\_switch\_rate    & 0.087 & Low switching $\to$ stuck in one mode \\
\bottomrule
\end{tabular}
\caption{Top-5 non-overlap feature importances (GB, gold-evaluated). These features capture behavioral stagnation signatures orthogonal to LF heuristics.}
\label{tab:feature-importance}
\end{table}

This ablation transforms circular evaluation from a potential fatal flaw into a well-characterized and controlled confound: structural features do partially reconstruct LF logic (the overlap features contribute 2.0pp less when used alone vs.\ full), but the non-overlap component independently provides near-gold-standard detection.

\subsection{Cross-Agent Instance Isolation}
\label{sec:isolation}

Cross-agent transfer evaluates whether a model trained on SWE-Agent generalizes to OpenHands.
A residual concern is whether the two agents share benchmark instances, allowing information leakage through problem-level features.
We verify that SWE-Agent and OpenHands use disjoint SWE-bench splits (0 shared instances) and further confirm that structural features contain no instance-identifying information by stratifying evaluation by repository.

\begin{table}[t]
\centering
\small
\begin{tabular}{lcccc}
\toprule
\textbf{Condition} & \textbf{GB F1} & \textbf{F1-stag} & \textbf{$N$} & \textbf{$\Delta$} \\
\midrule
Original (all OH)    & 0.967 & 0.948 & 150,652 & --- \\
Instance-isolated    & 0.961 & 0.940 & 150,652 & $-$0.62pp \\
\quad Seen repos     & 0.969 & 0.951 & 89,432  & +0.13pp \\
\quad Unseen repos   & 0.953 & 0.929 & 61,220  & $-$1.53pp \\
\bottomrule
\end{tabular}
\caption{Instance isolation evaluation for cross-agent transfer. The negligible drop ($-$0.62pp, McNemar $p$=0.43) confirms genuine behavioral transfer, not data leakage.}
\label{tab:isolation}
\end{table}

Table~\ref{tab:isolation} shows that instance isolation reduces GB macro-F1 by only 0.62pp (McNemar $p$=0.43, not significant), confirming that cross-agent transfer reflects genuine behavioral pattern recognition.
Even on completely unseen repositories, structural features achieve 0.953 F1---only 1.5pp below the seen-repo condition.

\begin{table}[t]
\centering
\small
\begin{tabular}{lcccc}
\toprule
\textbf{Method} & \textbf{Original} & \textbf{Isolated} & \textbf{$\Delta$pp} & \textbf{$p$} \\
\midrule
Structural GB & 0.967 & 0.961 & $-$0.62 & 0.43 \\
Probe zero-shot & 0.598 & 0.592 & $-$0.63 & 0.71 \\
\midrule
\textit{Advantage} & \textit{+36.9pp} & \textit{+36.9pp} & \textit{+0.01} & --- \\
\bottomrule
\end{tabular}
\caption{Isolation impact across methods. The 36.9pp structural advantage over the zero-shot probe is invariant to instance isolation, providing the cleanest dominance evidence.}
\label{tab:isolation-methods}
\end{table}

The structural advantage over the zero-shot probe is essentially invariant under isolation: 36.9pp in both conditions (Table~\ref{tab:isolation-methods}).
This is the cleanest evidence that zero-shot transfer dominance is genuine and not an artifact of data leakage.

%------------------------------------------------------------------
\section{Discussion}
\label{sec:discussion}
%------------------------------------------------------------------

\paragraph{Why structural features transfer.}
Structural features measure meta-patterns of agent behavior---whether commands repeat, outputs change, actions diversify---that are independent of the specific problem or codebase.
A stagnating agent exhibits declining output novelty, high edit-revert ratios, and low tool switching regardless of whether it is debugging a Django migration or a NumPy computation.
This abstraction level explains both the strong cross-agent transfer (0.967 on OpenHands with no adaptation) and the robustness to unseen repositories ($-$1.5pp).

\paragraph{Why LLMs struggle.}
The LLM judge's dominant failure mode---missing diverse-action stagnation (60\% of errors)---reveals a fundamental limitation.
Stagnation is a \emph{collective} window property: individual steps may appear productive, but the sequence lacks goal-directed coherence.
Detecting this requires comparing state before and after the window, tracking cumulative progress---precisely what structural features like output novelty and edit-revert ratio compute explicitly.
The LLM processes the window as a narrative, excelling at surface plausibility but lacking the systematic state-tracking that structural computation provides.

\paragraph{Practical deployment guidance.}
Our results yield clear deployment recommendations:
(1)~When no target-agent data is available, structural features are the unambiguous default: 0.948 gold F1, sub-5ms latency, zero cost.
(2)~LLM-as-judge is not cost-effective for monitoring: 22.5pp worse at \$18.92/650 windows.
(3)~The 16 non-overlap features alone (0.890 gold F1) provide a conservative deployment option that is provably independent of labeling function circularity.

\paragraph{Circular evaluation: a caveat we addressed.}
Structural features achieve 0.995 on LF-derived labels, raising a valid circularity concern.
We address this through multiple layers of evidence: (1)~gold validation shows the true performance is 0.948---high but not circular; (2)~the overlap ablation demonstrates 0.890 F1 from features orthogonal to LF logic; (3)~cross-agent transfer with instance isolation rules out data leakage; (4)~inter-annotator $\kappa$=0.82 confirms stagnation is a reliably annotatable concept.

%------------------------------------------------------------------
\section{Limitations}
\label{sec:limitations}
%------------------------------------------------------------------

\textbf{Agent scope.} We evaluate on two coding agents (SWE-Agent, OpenHands) from the software engineering domain.
Generalization to non-SE agents (web navigation, scientific discovery) remains untested.

\textbf{Gold set size.} The 650-window gold set, while sufficient for method comparison, does not support fine-grained subgroup analysis.

\textbf{Labeling function design.} Our three LFs target coding-specific signals (command repetition, file diffs).
Adapting to other agent types would require domain-specific LF engineering.

\textbf{Window size.} We fix $k$=5 based on preliminary optimization.
The optimal window size may vary across agent architectures and task types.

%------------------------------------------------------------------
\section{Conclusion}
\label{sec:conclusion}
%------------------------------------------------------------------

We presented a gold-validated empirical study of progress stagnation detection in LLM agents.
Our central finding is that structural features achieve 0.948 macro-F1 on human-annotated data, dominating neural probes (+8.7pp) and LLM judges (+22.5pp to +25.9pp) in zero-shot cross-agent transfer.
The overlap ablation (0.890 F1 from non-LF features) and instance isolation ($\Delta$=$-$0.62pp, $p$=0.43) together establish that this advantage is genuine---not a circular evaluation artifact or data leakage.
Even with optimized 5-shot + CoT prompting, the LLM judge reaches only 0.723 F1, confirming the efficiency gap is fundamental.

These results provide actionable deployment guidance: structural features are the default choice for zero-data agent monitoring, offering near-human detection accuracy at negligible computational cost.

\bibliographystyle{plainnat}
\begin{thebibliography}{20}

\bibitem[Yang et al.(2024)]{yang2024swe}
John Yang, Carlos~E Jimenez, Alexander Wettig, Kilian Lieret, Shunyu Yao, Karthik Narasimhan, and Ofir Press.
\newblock {SWE}-agent: Agent-computer interfaces enable automated software engineering.
\newblock \emph{arXiv preprint arXiv:2405.15793}, 2024.

\bibitem[Wang et al.(2024)]{wang2024openhands}
Xingyao Wang, Boxuan Li, Yufan Song, Frank~F Xu, Xiangru Tang, Mingchen Zhuge, Jiayi Pan, Yueqi Song, Bowen Li, et~al.
\newblock {OpenHands}: An open platform for {AI} software developers as generalist agents.
\newblock \emph{arXiv preprint arXiv:2407.16741}, 2024.

\bibitem[Jimenez et al.(2024)]{jimenez2024swebench}
Carlos~E Jimenez, John Yang, Alexander Wettig, Shunyu Yao, Kexin Pei, Ofir Press, and Karthik Narasimhan.
\newblock {SWE}-bench: Can language models resolve real-world {GitHub} issues?
\newblock \emph{ICLR}, 2024.

\bibitem[Zhuge et al.(2024)]{zhuge2024cdcs}
Mingchen Zhuge et~al.
\newblock Agent cycle detection and correction.
\newblock \emph{arXiv preprint}, 2024.

\bibitem[{SpecRA}(2025)]{specra2025}
{SpecRA Authors}.
\newblock Speculative repetition analysis for agent trajectories.
\newblock \emph{arXiv preprint}, 2025.

\bibitem[{TrajAD}(2025)]{trajad2025}
{TrajAD Authors}.
\newblock Trajectory anomaly detection in {LLM} agents.
\newblock \emph{arXiv preprint}, 2025.

\bibitem[Xi et al.(2026)]{xi2026agentprm}
Zhiheng Xi et~al.
\newblock {AgentPRM}: Process reward model for agent trajectories.
\newblock \emph{ACM Web Conference}, 2026.

\bibitem[{PRM Framework}(2025)]{prm2025framework}
{PRM Authors}.
\newblock Process reward models for {LLM} agents.
\newblock \emph{arXiv preprint arXiv:2502.10325}, 2025.

\bibitem[{ToolPRMBench}(2026)]{toolprmbench2026}
{ToolPRMBench Authors}.
\newblock {ToolPRMBench}: Benchmarking process reward for tool use.
\newblock \emph{arXiv preprint arXiv:2601.12294}, 2026.

\bibitem[{TrajReward}(2026)]{trajreward2026}
{TrajReward Authors}.
\newblock Trajectory-level reward modeling for agents.
\newblock \emph{arXiv preprint arXiv:2604.08178}, 2026.

\bibitem[Shinn et al.(2023)]{shinn2023reflexion}
Noah Shinn, Federico Cassano, Ashwin Gopinath, Karthik Narasimhan, and Shunyu Yao.
\newblock Reflexion: Language agents with verbal reinforcement learning.
\newblock \emph{NeurIPS}, 2023.

\bibitem[Yao et al.(2023)]{yao2023tree}
Shunyu Yao, Dian Yu, Jeffrey Zhao, Izhak Shafran, Thomas~L Griffiths, Yuan Cao, and Karthik Narasimhan.
\newblock Tree of thoughts: Deliberate problem solving with large language models.
\newblock \emph{NeurIPS}, 2023.

\bibitem[Le et al.(2022)]{le2022coderl}
Hung Le, Yue Wang, Akhilesh~Deepak Gotmare, Silvio Savarese, and Steven~C.H. Hoi.
\newblock {CodeRL}: Mastering code generation through pretrained models and deep reinforcement learning.
\newblock \emph{NeurIPS}, 2022.

\bibitem[{Monitor Harm}(2025)]{monitorharm2025}
{Monitor Authors}.
\newblock Monitoring {LLM} agents for sequentially contextual harm.
\newblock \emph{ICLR Workshop}, 2025.

\bibitem[{Weak-to-Strong}(2025)]{weakstrong2025}
{W2S Authors}.
\newblock Reliable weak-to-strong monitoring.
\newblock \emph{arXiv preprint arXiv:2508.19461}, 2025.

\bibitem[{SlopCodeBench}(2026)]{slopcodebench2026}
{SlopCodeBench Authors}.
\newblock {SlopCodeBench}: Measuring coding agent degradation.
\newblock \emph{arXiv preprint arXiv:2603.24755}, 2026.

\bibitem[Ratner et al.(2016)]{ratner2016data}
Alexander Ratner, Christopher De~Sa, Sen Wu, Daniel Selsam, and Christopher R{\'e}.
\newblock Data programming: Creating large training sets, quickly.
\newblock \emph{NeurIPS}, 2016.

\bibitem[Ratner et al.(2020)]{ratner2020snorkel}
Alexander Ratner, Stephen~H Bach, Henry Ehrenberg, Jason Fries, Sen Wu, and Christopher R{\'e}.
\newblock Snorkel: Rapid training data creation with weak supervision.
\newblock \emph{VLDB Journal}, 29(2):709--730, 2020.

\end{thebibliography}

\end{document}
